\section{Acquiring control changes the measurement and the period}
\label{sim:acquisition-period-design}

Suppose the financing arrangement eventually gives the creditor control of the debtor. The debtor's legal accounts do not thereby become the creditor's accounts, and the new shares do not supply an acquisition-date balance sheet. Consolidated reporting now needs to explain which assets and liabilities enter the group, at what amounts, and from which instant the acquired company's income belongs in the group result. Multiplying its annual statement by a 75 per cent shareholding would answer none of those questions. It would both omit part of the controlled company's activity and bring in activity from before control.

The smallest continuation considered here is one gain of control in an ordinary business acquisition, with a single reporting currency, aligned reporting dates and no subsequent loss of control or change in ownership rights during the selected interval. These are limits of the first method. A purchase of assets that is not a business, a combination under common control, an investment-entity exception and separate financial statements require their own treatments. The business assessment and acquirer identification therefore precede the acquisition calculation. The acquisition date follows the assessed date of obtaining control, rather than the share-issue date by definition \citep[paras.~2--9]{aasb3scope2026}.

\subsection{The acquisition balance needs its own measurements}

Consider the consolidated acquisition worksheet at that date. Let $A$ be the aggregate recognised acquired assets and $L$ the aggregate assumed liabilities, both expressed as positive magnitudes in the same currency unit. Their amounts follow the applicable recognition and measurement rules, including item-specific exceptions; an unchanged legal carrying amount is not automatically the required acquisition amount. Let $T$ denote measured consideration transferred, $Q$ measured non-controlling interests, and $P$ the acquisition-date fair value of a qualifying previously held equity interest. Under the ordinary business-acquisition branch, the excess described in AASB~3 paragraph~32 is
\begin{equation}
 E=T+Q+P-(A-L).
 \label{eq:acquisition-conditional-excess}
\end{equation}
The calculation compares the consideration and interests in the acquired business with its recognised net assets. A positive excess can determine goodwill only when the route and measurements support that recognition. A negative excess requires the reassessment specified in paragraphs~34--36 before a bargain-purchase gain is recognised. Computing $\max(E,0)$ and putting the remainder into profit would skip that requirement \citep[paras.~10--21, 24--36]{aasb3scope2026}.

The distinction between $T$ and $P$ matters particularly in the conversion example. A creditor may hold debt immediately before obtaining shares. That debt is not, merely because it converts, the previously held equity interest addressed by paragraph~42. Nor does the issuer's transfer of a carrying amount from debt to equity determine the holder's acquisition consideration. The contract and its financial-instrument treatment must establish what is transferred and which effects belong to the business combination. Pre-existing relationships and separate transactions require their own assessment \citep[paras.~42, 51--53 and B50--B53]{aasb3scope2026}. The face amount, an economic price and a carrying amount cannot silently fill the same measurement slot.

Non-controlling interests introduce a further measurement choice. For qualifying present ownership interests, a proportional-share election gives $Q=q(A-L)$, where $q$ is the non-controlling fraction. A fair-value election instead requires a distinct measurement of $Q$ \citep[para.~19]{aasb3scope2026}. Substituting the proportional expression into equation~\eqref{eq:acquisition-conditional-excess} gives
\begin{equation}
 E=T+P-(1-q)(A-L).
 \label{eq:acquisition-proportional-excess}
\end{equation}
For a conditional 75 per cent ordinary interest, $q=1/4$ and the acquired net-asset contribution is $3(A-L)/4$. This algebra does not establish that a particular preference share or other right qualifies for the election. With full ownership, $q=0$ and no non-controlling interest remains. An unspecified prior interest does not become $P=0$; that zero requires an explicit no-prior-equity premise under the selected route.

\subsection{Only the controlled interval contributes subsidiary results}

After the acquisition worksheet is established, consolidation incorporates the controlled company's income and expenses from the date control is obtained. Those subsequent amounts follow the acquisition-date basis: for example, depreciation uses the acquired asset amount recognised in the consolidated statements \citep[para.~20 and B86--B88]{aasb10control2026}. A legal-book expense may therefore need a group adjustment even when the legal journal itself is correct. The simulator must preserve the legal books and put the adjustment in the separate consolidation worksheet.

Write $R$ for the subsidiary's legal profit over the exact post-control interval, and $D$ for the signed additional expense required to express that interval on the group basis. Positive $D$ reduces profit; negative $D$ increases it. Let $O$ be the legal other comprehensive income over the same interval and $U$ its signed group adjustment. All four quantities have the same currency unit. The adjusted subsidiary contribution is
\begin{equation}
 \pi=R-D,\qquad \omega=O+U.
 \label{eq:acquisition-post-control-results}
\end{equation}
This is a period and measurement reconciliation, not a forecast or a rule for dividing an annual total. A whole-year $R$ cannot be prorated by elapsed months unless an independently supported method justifies that approximation. A journal recorded only by calendar date also needs an ordering premise when control changes within that date; rounding the control instant would conceal the ambiguity.

Under the deliberately limited assumption of constant ordinary present ownership rights, the owners' and non-controlling shares of $\pi$ are $(1-q)\pi$ and $q\pi$, and similarly for $\omega$. Actual allocation must consider the relevant rights and the potential-voting-rights exception; the fraction of issued shares is not a universal allocation rule \citep[B89--B95]{aasb10control2026}. Let $V$ be distributions to non-controlling interests, expressed as a positive reduction, and $M$ the other signed, supported changes in those interests within this limited branch. Their selected reconciliation is then
\begin{equation}
 N_{\mathrm{close}}=Q+q(R-D)+q(O+U)-V+M.
 \label{eq:acquisition-nci-reconciliation}
\end{equation}
Starting with measured acquisition-date NCI, the equation adds its share of adjusted profit and OCI and then accounts for direct changes. A sufficiently large loss may make the closing amount negative; clipping that deficit at zero would contradict the attribution requirement \citep[B94]{aasb10control2026}. The equation cannot establish whether the supplied $M$ or adjustment population is complete. Ownership transactions with continued control, and a later loss of control, need their separate methods rather than an unexplained entry in $M$.

\subsection{Cash and revisions retain their own dates}

Acquisition also changes the cash included in the group. The conversion of debt to equity is a noncash transaction, but the acquired company may already hold cash. Its cash enters the consolidated perimeter even when the conversion itself creates no receipt or payment. The method must retain cash consideration, acquired cash and noncash consideration separately, with the applicable acquisition cash presentation net of acquired cash. Acquisitions and disposals are not netted against one another merely because both change control \citep[paras.~39--44B]{aasb107acquisition2026}. Cash transfers before control are not retrospectively intragroup flows; transfers afterwards require supported counterpart matching before elimination.

The acquisition worksheet must also distinguish new information about conditions at the acquisition date from later events. Provisional amounts and measurement-period adjustments follow their own rules, including the one-year maximum, rather than a general permission to rewrite the past \citep[paras.~45--50]{aasb3scope2026}. A revised measurement can therefore have an acquisition-date effect and a later information-availability date. The earlier actor view must remain reproducible, using the temporal distinctions developed in Chapter~\ref{sim:continuous-methods}.

\subsection{A transaction date does not select an edition}

The conditional AASB analysis preserves several dates because they answer different questions. For example, the retained AASB~2026-2 compilation was prepared on 21 October 2026, has a compilation date of 31 December 2027, and generally applies to annual reporting periods beginning on or after 1 July 2029, with earlier application permitted under its terms. These dates cannot be substituted for one another. Its amendments to AASB~3 and AASB~10 concern application of the Conceptual Framework; entity classification and elections remain necessary inputs \citep[paras.~3, 18--19, 25; compilation details]{aasb2026frameworkcompiled}. The inspected interaction with AASB~2026-4 also imposes a specific AASB~18 condition for early use of that amendment \citep[paras.~3, 25]{aasb2026presentation}.

Other branches remain visible. AASB~20's Appendix~D adds a business-combination treatment for regulatory assets and liabilities, with its own application timing and an additional dependency when applied before AASB~18. The retained appendix alone does not establish the entity's eligibility or the complete transition treatment \citep[Appendix~D, amendments to AASB~3]{aasb20acquisitionappendix}. AASB~1061's Appendix~C changes tier wording relevant to AASB~10; it does not supply the full Tier~3 scope or measurement rules \citep[Appendix~C, amendment to AG1]{aasb1061scopeappendix}. Similarly, the inspected sale/contribution amendment chain affects a particular loss-of-control branch, rather than replacing every acquisition method \citep[paras.~6--8]{aasb2024controltiming}. The dated authority matrix retains these alternatives and exact paragraph anchors. A future transaction's policy must be checked against its entity, annual period, elections and then-current applicable requirements.

\subsection{What the selected reference establishes}

The conditional expression reference binds the original control-history request and the exact gain-of-control instant before constructing these relations. Each amount remains an explicitly named unknown, with its role, currency, unit, timing, measurement basis and source identity. Future or unknown source availability remains unresolved for the selected actor. Missing assets, consideration, tax effects or adjustment methods cannot be supplied by the fact that an equation balances.

Independent coefficient maps check both NCI elections, qualifying prior equity, the full-ownership limit and the post-control reconciliations. Adverse requests test prior debt labelled as equity, pre-control profit leakage, sub-day rounding, inconsistent currency or ownership, source alteration and attempted insertion of an unmeasured face value. All source, accounting, posting and report authorization remains withheld.\footnote{The dated source review is retained in \path{data/accounting-master/acquisition-period-authority-20261026/}; its 13 primary PDFs have retained identities and inspected anchors. The symbolic design and independent expected expressions are in \path{data/accounting-master/acquisition-symbolic-20261027/}; the selected check covers 25 cases and 21 refusal controls. These deterministic checks establish the stated local behaviour, not source truth or complete acquisition accounting.}

The next financial-method comparison must expand the aggregate unknowns into an identified item population, preserve every legal account, and reconcile a complete dated acquisition worksheet and period report against independently authored expectations. That comparison must include consideration and separate transactions, acquisition-basis expenses, NCI, cash, eliminations, taxes and revisions. The expressions developed here identify the measurements that comparison needs; they do not replace it.
